\section{总结与延伸}

\subsection{核心结论}

本节把全篇的机器人判断压缩成检查表。读者可以用它来评估一个机器人 demo 或公司叙事：它解决的是大脑、小脑、数据、本体、触觉、安全，还是只展示了其中一块？

\begin{enumerate}
\item 机器人过去十年的第一范式转移，是强化学习和 Sim2Real 改变 locomotion。
\item 大模型带来第二范式转移，让机器人获得语言理解、常识和计划能力。
\item 机器人基座模型目前仍主要是多模态大模型加动作数据和 fine-tuning，尚未完全独立。
\item 高质量数据是机器人当前最重要瓶颈，几百小时数据不足以验证 scaling。
\item Gemini Robotics 1.5 的 Motion Transfer 试图解决跨本体泛化。
\item Synthetic Data 是谭杰认为必须相信的关键方向，光靠真实数据不够。
\item 世界模型的价值在于预测动作后果，而不只是生成视频。
\item 灵巧手让触觉重新重要，因为精细 manipulation 需要接触信息。
\item 通用机器人会经历从固定自动化到超人能力的长期阶段。
\item 中美机器人最理想的路径是智能范式和硬件供应链互补。
\end{enumerate}

\subsection{开放问题}

本节保留开放问题，是因为本期讨论的许多技术路线仍处在快速变化阶段。机器人行业最需要的不是单个答案，而是持续跟踪哪些变量真的改善：数据规模、跨本体迁移、触觉硬件、世界模型和安全评测。

\begin{itemize}
\item 机器人基座模型什么时候会真正独立于多模态大模型？
\item Motion Transfer 能否在足够多本体之间稳定工作？
\item Synthetic Data 如何缩小 sim-to-real gap？
\item 世界模型是视觉预测更重要，还是动作/状态预测更重要？
\item 触觉硬件何时能规模化进入灵巧手数据飞轮？
\item 通用人形机器人和垂直 specialist 会并存多久？
\item 机器人 safety 应该如何在能力快速增长前先行？
\end{itemize}

\subsection{拓展阅读}

\begin{itemize}
\item Sim2Real Learning Agile Locomotion for Quadruped Robots：理解强化学习如何迁移到真实四足机器人。
\item Robot Transformer 系列：RT-1、RT-2、RT-X 等工作。
\item Gemini Robotics / Gemini Robotics 1.5：Google DeepMind 机器人模型路线。
\item EP132 高继扬访谈：星海图、Waymo/Momenta 和具身智能生产化。
\item EP134 数据综述：机器人数据、Recipe 和数据定价。
\end{itemize}

\begin{importantbox}{最后的判断}
机器人最难的不是“让模型看懂一句话”，而是让一个实体系统在真实世界中稳定行动。大模型给了机器人大脑，强化学习给了部分小脑，合成数据和世界模型试图补足经验，但真正落地还需要高质量数据、可靠硬件、触觉、安全和长期耐心。
\end{importantbox}

\end{document}
